% !TEX root = ../../main.tex

\section{本章小结}

本章围绕 HearBridge 系统进行了需求分析。
首先，从系统总体角度明确了普通用户和系统管理员两类主要角色，并说明了系统面向手语训练、句子视频识别、资源管理和模型版本维护等场景的功能范围；
随后分别分析了普通用户和管理员的功能需求，并通过用例图说明两类角色与系统功能之间的关系；
接着从用户视角描述了普通用户总体操作流程、实时手势训练流程和句子视频识别流程；
随后选取用户登录、实时手势训练、句子视频识别、手势资源管理和模型版本发布等关键用例进行规约说明；
最后，从性能、可靠性、易用性、安全性、可维护性与扩展性等方面分析了系统的非功能性需求，并对当前系统实现边界进行了约束说明。

通过本章分析，可以明确 HearBridge 系统需要同时满足普通用户训练识别体验和管理员资源模型维护两方面需求。
下一章将在需求分析的基础上，对系统总体架构、功能模块、数据结构和关键业务流程进行设计。
